% ==============================================================================
% 01_supervision.tex — UE SIS589
% ==============================================================================
\chapter{Supervision des Systèmes d'Information}
\label{chap:supervision}

\epigraph{\itshape « On ne peut défendre ce qu'on ne voit pas. »}
         {--- Principe fondateur de la supervision active}

\begin{boxobjectifs}
\begin{enumerate}
  \item Distinguer supervision, monitoring et surveillance~; positionner chaque
        concept dans le cycle de sécurité opérationnelle.
  \item Décrire l'architecture fonctionnelle d'un SOC et ses modèles
        organisationnels.
  \item Identifier les sources de données d'un SIEM et comprendre le pipeline
        de traitement événement~$\to$~alerte.
  \item Définir et calculer les métriques fondamentales~: MTTD, MTTR, FPR,
        couverture MITRE ATT\&CK.
  \item Appliquer la grille de maturité SOC à LiZbiZ-SIS.
\end{enumerate}
\end{boxobjectifs}

% ==============================================================================
\section{Fondements conceptuels}
\label{sec:superv-fondements}
% ==============================================================================

\subsection{Définitions structurantes}

\begin{boxdef}
\textbf{Supervision de la sécurité~:} Ensemble des processus, outils et pratiques
permettant d'observer en continu l'état du SI, de détecter toute déviation par
rapport à une \textit{baseline} de comportement normal, et de déclencher une réponse
dans un délai compatible avec les objectifs de l'organisation.
\end{boxdef}

\begin{boxdef}
\textbf{Événement de sécurité (\textit{Security Event})~:} Changement observable
d'état dans un composant du SI susceptible d'avoir une incidence sur la sécurité.
Un événement n'est pas, par définition, un incident.
\end{boxdef}

\begin{boxdef}
\textbf{Incident de sécurité~:} Événement ou série d'événements confirmés portant
atteinte — ou susceptibles de porter atteinte — à la confidentialité, l'intégrité ou
la disponibilité du SI (définition ISO~27035). Tout incident est un événement~;
tout événement n'est pas un incident.
\end{boxdef}

La supervision opère à trois niveaux~:

\begin{enumerate}
  \item \textbf{Disponibilité~:} ping, healthcheck, état des processus.
        Outils~: Nagios, Zabbix, Prometheus.
  \item \textbf{Performance~:} CPU, mémoire, bande passante, temps de réponse.
        Outils~: Grafana, Datadog, Elastic APM.
  \item \textbf{Sécurité~:} analyse des journaux, corrélation, détection
        comportementale. Outil central~: \textbf{SIEM}. Objet de ce cours.
\end{enumerate}

\subsection{Le paradigme \textit{assume breach}}

\begin{boxdef}
\textit{\textbf{Assume breach}~:} Posture de sécurité consistant à considérer
que l'attaquant est potentiellement déjà présent dans le SI. L'objectif ne
se limite plus à empêcher l'intrusion mais à \textbf{détecter et contenir}
avant que l'attaquant n'atteigne ses objectifs.
\end{boxdef}

Ce paradigme impose trois conséquences opérationnelles~:
\begin{itemize}
  \item La surveillance interne (est-ouest) devient aussi critique que
        la surveillance périmétrique (nord-sud).
  \item Le \termecle{temps de séjour} (\textit{dwell time}) — durée entre
        compromission initiale et détection — devient l'indicateur de
        performance dominant.
  \item Le \textit{threat hunting} proactif complète la détection réactive
        par alertes.
\end{itemize}

\begin{boiteRemarque}{Temps de séjour~: chiffres de référence}
Mandiant M-Trends 2024~: médiane du \textit{dwell time} mondiale à
\textbf{10~jours} (amélioration significative vs 21~jours en 2023,
101~jours en 2013). En Afrique subsaharienne, les estimations sectorielles
restent proches de 30~jours faute de capacités de détection déployées.
La réduction de ce délai est l'objectif mesurable n°~1 de tout SOC.
\end{boiteRemarque}

% ==============================================================================
\section{Le Security Operations Center (SOC)}
\label{sec:soc}
% ==============================================================================

\subsection{Définition et mission}

\begin{boxdef}
\textbf{SOC (\textit{Security Operations Center})~:} Structure organisationnelle
— équipe humaine, processus et technologie — dédiée à la surveillance continue du
SI, à la détection des menaces, à l'analyse des alertes et à la coordination des
réponses aux incidents. Disponibilité cible~: 24h/24, 7j/7.
\end{boxdef}

Le SOC n'est pas un outil~: c'est une \textbf{capacité opérationnelle}. Son
efficacité repose sur trois piliers interdépendants~:
\begin{itemize}
  \item \textbf{Personnes~:} analystes N1/N2/N3, ingénieurs détection, manager SOC.
  \item \textbf{Processus~:} procédures, playbooks, escalades, SLA.
  \item \textbf{Technologie~:} SIEM, EDR, threat intelligence, SOAR.
\end{itemize}

\subsection{Niveaux d'analyse}

\begin{table}[H]
\centering\small
\caption{Niveaux d'analyse dans un SOC}
\label{tab:soc-tiers}
\rowcolors{2}{TableauPair}{TableauImpair}
\begin{tabular}{p{1.3cm}p{3.2cm}p{6.5cm}p{2.5cm}}
\toprule
\tableauHeader{Niveau} & \tableauHeader{Rôle} &
\tableauHeader{Activités principales} & \tableauHeader{Outillage} \\
\midrule
\textbf{N1} & Analyste de triage &
Surveillance tableaux de bord, qualification alertes, escalade N2,
création et documentation des tickets &
SIEM, ticketing \\
\textbf{N2} & Analyste d'investigation &
Analyse approfondie des alertes escaladées, investigation incidents
confirmés, corrélation avancée, recommandation confinement &
SIEM, EDR, sandbox \\
\textbf{N3} & Expert / Threat Hunter &
Investigation forensique, threat hunting proactif, développement règles
de détection, veille CTI &
Outils forensiques, CTI, SOAR \\
\textbf{Manager} & Direction opérationnelle &
Pilotage KPI, liaison RSSI/DG, gestion escalades critiques, reporting &
Tableaux de bord \\
\bottomrule
\end{tabular}
\end{table}

\subsection{Modèles organisationnels}

\begin{description}
  \item[\termecle{SOC interne}] Équipe dédiée, opérée par l'organisation.
        Avantage~: connaissance métier profonde, contrôle total.
        Inconvénient~: coût élevé (recrutement, rétention), couverture 24/7 difficile.
  \item[\termecle{SOC externalisé (MSSP)}] \textit{Managed Security Service Provider}.
        Avantage~: disponibilité immédiate, mutualisation des ressources.
        Inconvénient~: moindre connaissance métier, dépendance contractuelle.
  \item[\termecle{SOC hybride}] Équipe interne pour la connaissance métier et les
        escalades~; MSSP pour la surveillance nocturne et week-end. Modèle dominant
        dans les organisations de taille intermédiaire.
  \item[\termecle{Virtual SOC (vSOC)}] Équipe distribuée géographiquement,
        sans salle physique dédiée. Rendu viable par les plateformes SIEM cloud.
\end{description}

\subsection{Architecture technique}

\begin{figure}[H]
\centering
\begin{tikzpicture}[
  box/.style={rectangle, rounded corners=3pt, draw=CouleurPrimaire,
              fill=CouleurDef, minimum width=2.9cm, minimum height=0.75cm,
              font=\small\bfseries, text centered},
  arrow/.style={-Stealth, thick, color=CouleurSecondaire}]

\node[box,fill=CouleurLizbiz,draw=CouleurBordLizbiz] (fw)  at (-5, 2.2) {Firewalls/IDS};
\node[box,fill=CouleurLizbiz,draw=CouleurBordLizbiz] (srv) at (-5, 0.7) {Serveurs OS};
\node[box,fill=CouleurLizbiz,draw=CouleurBordLizbiz] (app) at (-5,-0.8) {Applications};
\node[box,fill=CouleurLizbiz,draw=CouleurBordLizbiz] (edr) at (-5,-2.3) {EDR / AV};

\node[box,minimum width=2.4cm,fill=CouleurSiem,draw=CouleurBordSiem]
  (col) at (-1.2,-0.2) {Collecteur\\(Agent/Syslog)};

\node[box,minimum width=3.8cm,minimum height=1.3cm,
      fill=CouleurPrimaire,draw=CouleurPrimaire,text=white]
  (siem) at (2.8,-0.2) {\Large\textbf{SIEM}};

\node[box,fill=CouleurAlerte,draw=CouleurBordAlerte] (alert)  at (6.5, 1.2) {Alertes};
\node[box,fill=CouleurInfo,  draw=CouleurBordInfo]   (dash)   at (6.5,-0.2) {Dashboards};
\node[box,fill=CouleurInfo,  draw=CouleurBordInfo]   (report) at (6.5,-1.6) {Rapports};
\node[box,fill=CouleurForensic,draw=CouleurBordForensic,minimum width=2.4cm]
  (soar) at (2.8, 2.2) {SOAR / Ticketing};

\foreach \s in {fw,srv,app,edr} \draw[arrow] (\s) -- (col);
\draw[arrow] (col)  -- (siem);
\draw[arrow] (siem) -- (alert);
\draw[arrow] (siem) -- (dash);
\draw[arrow] (siem) -- (report);
\draw[arrow] (siem) -- (soar);

\node[font=\tiny\itshape,color=CouleurBordInfo] at (-3.0,-1.0) {Normalisation ECS/CEF};
\node[font=\tiny\itshape,color=CouleurBordInfo] at (4.8, 0.5) {Corrélation + enrichissement};
\end{tikzpicture}
\caption{Architecture type d'un SOC~: de la collecte à l'alerte}
\label{fig:soc-archi}
\end{figure}

\begin{boxlizbiz}
\textbf{LiZbiZ-SIS~: SOC minimal cible}

Contraintes~: budget $<$ 5~M~FCFA/an, 2~techniciens disponibles,
couverture heures ouvrables + astreinte. Solution retenue~:
\textbf{Wazuh~Manager + ELK Stack} (open-source)~; agents sur les
serveurs Windows (AD, ERP) et Linux (Apache, MySQL)~; Kibana pour la
visualisation~; TheHive pour le ticketing~; OpenCTI pour la CTI.
Objectif~: niveau~3 de maturité (tableau~\ref{tab:soc-maturite})
à l'issue des laboratoires.
\end{boxlizbiz}

% ==============================================================================
\section{Le SIEM~: cœur technique}
\label{sec:siem}
% ==============================================================================

\subsection{Définition et pipeline}

\begin{boxdef}
\textbf{SIEM (\textit{Security Information and Event Management})~:}
Plateforme centralisant la collecte de journaux issus de sources hétérogènes,
les normalisant dans un format commun, les corrélant selon des règles
ou des modèles comportementaux, et générant des alertes lorsqu'un scénario
de menace est détecté.
\end{boxdef}

Le pipeline de traitement en six étapes~:

\begin{enumerate}
  \item \textbf{Collecte.} Syslog UDP/TCP~514, agent propriétaire, API REST,
        fichier plat.
  \item \textbf{Parsing.} Décomposition de la ligne brute en champs structurés
        (timestamp, IP source, utilisateur, action, statut).
  \item \textbf{Normalisation.} Conversion vers un schéma commun~: ECS
        (Elastic Common Schema), CEF (ArcSight), LEEF (QRadar). Horodatage UTC.
  \item \textbf{Enrichissement.} Géolocalisation IP, résolution DNS inverse,
        lookup CTI (IoC connus), classe de l'actif dans le CMDB.
  \item \textbf{Corrélation.} Application des règles de détection sur des fenêtres
        temporelles~: \og{}5~échecs d'authentification depuis la même IP
        en 60~secondes\fg{}.
  \item \textbf{Alerte.} Création de ticket, notification analyste N1,
        déclenchement optionnel d'une action SOAR.
\end{enumerate}

\begin{lstlisting}[style=StyleTerminal,
  caption={Log brut Apache $\to$ événement normalisé ECS $\to$ enrichi CTI},
  label={lst:siem-pipeline}]
# 1. Ligne brute (Apache Combined Log)
203.0.113.42 - - [12/Apr/2026:03:14:22 +0000] \
  "POST /wp-login.php HTTP/1.1" 401 1234 \
  "-" "sqlmap/1.7 (https://sqlmap.org)"

# 2. Après parsing + normalisation ECS (Logstash grok)
{
  "@timestamp"                  : "2026-04-12T03:14:22Z",
  "source.ip"                   : "203.0.113.42",
  "url.path"                    : "/wp-login.php",
  "http.request.method"         : "POST",
  "http.response.status_code"   : 401,
  "event.outcome"               : "failure",
  "event.category"              : "authentication",
  "user_agent.original"         : "sqlmap/1.7"
}

# 3. Après enrichissement CTI (lookup AlienVault OTX)
{
  "source.geo.country_iso_code" : "RU",
  "threat.indicator.type"       : "ip",
  "threat.feed.name"            : "AlienVault OTX",
  "threat.indicator.confidence" : 85,
  "tags"                        : ["sqlmap","brute_force_candidate"]
}
\end{lstlisting}

\subsection{Sources d'événements prioritaires}

\begin{table}[H]
\centering\small
\caption{Sources SIEM et menaces détectables}
\label{tab:sources-siem}
\rowcolors{2}{TableauPair}{TableauImpair}
\begin{tabular}{p{3cm}p{3.2cm}p{7.3cm}}
\toprule
\tableauHeader{Source} & \tableauHeader{Format} &
\tableauHeader{Menaces détectables} \\
\midrule
Firewall / UTM    & Syslog CEF     & Flux bloqués, scan de ports, exfiltration \\
Serveur Windows   & WEF, Event IDs & Brute force (4625), création compte (4720),
                                     élévation (4672), effacement log (1102) \\
Serveur Linux     & Syslog, auth.log & SSH brute force, sudo anormal, cron suspect \\
Application web   & Apache/Nginx log & SQLi, scanner web, codes 4xx en masse \\
Active Directory  & LDAP/Kerberos  & Pass-the-hash, Kerberoasting, LDAP recon \\
EDR / AV          & Agent propre   & Malware, injection mémoire, process hollow \\
DNS               & Bind/Windows   & DGA, DNS tunneling, exfiltration \\
Proxy / DLP       & ICAP/Syslog    & Fuite de données, catégories interdites \\
\bottomrule
\end{tabular}
\end{table}

\subsection{Solutions SIEM de référence}

\begin{description}
  \item[\termecle{Splunk}] Leader commercial, langage SPL expressif.
        Coût proportionnel au volume ingéré (TB/jour).
  \item[\termecle{Elastic SIEM (ELK)}] Open-source. Elasticsearch + Logstash
        + Kibana + Elastic Security. Solution privilégiée pour les
        budgets contraints. Déployée pour LiZbiZ-SIS.
  \item[\termecle{IBM QRadar}] Robuste, présent dans les grandes entreprises
        et administrations africaines. Gestion native des \textit{offenses}.
  \item[\termecle{Microsoft Sentinel}] SIEM cloud-native Azure.
        Intégration native M365, Azure AD, Defender. Facturation à l'ingestion.
  \item[\termecle{Wazuh}] HIDS/SIEM open-source, agent unifié (FIM, compliance,
        détection d'intrusion). Déployé comme couche de collecte et de
        corrélation dans LiZbiZ-SIS.
\end{description}

% ==============================================================================
\section{Règles de corrélation}
\label{sec:regles-correlation}
% ==============================================================================

\subsection{Structure d'une règle}

Une règle de détection est définie par cinq attributs~:
\begin{enumerate}
  \item \textbf{Condition logique~:} combinaison d'événements (ET, OU, séquence,
        négation).
  \item \textbf{Fenêtre temporelle~:} durée dans laquelle les événements
        corrélés doivent survenir.
  \item \textbf{Seuil~:} nombre d'occurrences déclenchant l'alerte.
  \item \textbf{Regroupement~:} champ de corrélation (par IP, utilisateur,
        machine).
  \item \textbf{Gravité~:} critique / haute / moyenne / basse.
\end{enumerate}

\begin{lstlisting}[style=StyleCode, language=YAML,
  caption={Règle Wazuh --- brute force SSH, mapping MITRE ATT\&CK},
  label={lst:regle-ssh}]
<rule id="100200" level="10" frequency="5" timeframe="60">
  <if_matched_sid>5716</if_matched_sid>  <!-- echec auth SSH -->
  <same_source_ip />
  <description>Brute force SSH : 5 echecs en 60s meme IP</description>
  <mitre>
    <id>T1110.001</id>  <!-- Brute Force: Password Guessing -->
  </mitre>
  <group>authentication_failures,brute_force,pci_dss_10.2.4</group>
</rule>
\end{lstlisting}

\subsection{Faux positifs et fatigue d'alerte}

\begin{boxdef}
\textbf{Faux positif (FP)~:} Alerte signalant une menace inexistante.
Un taux de FP élevé ($>80\,\%$) génère une \termecle{fatigue d'alerte}
(\textit{alert fatigue})~: les analystes finissent par ignorer les alertes,
y compris les vraies.
\end{boxdef}

\begin{boxdef}
\textbf{Faux négatif (FN)~:} Menace réelle non détectée. Conséquence~:
compromission silencieuse et temps de séjour élevé. Plus dangereux que le FP
mais moins visible.
\end{boxdef}

La gestion du compromis FP/FN est au cœur de l'ingénierie de détection~:
abaisser un seuil réduit les FN mais augmente les FP. Chaque règle doit
être calibrée sur la \textit{baseline} de l'environnement cible.

\begin{boxsiem}
\textbf{Alerte SIEM --- LiZbiZ-SIS, 12/04/2026 03:14 UTC}
\begin{verbatim}
RULE    : 100200 — Brute force SSH
LEVEL   : 10 (HIGH)
SOURCE  : 203.0.113.42  (RU, TOR exit — AlienVault OTX)
TARGET  : srv-web-01.lizbiz.local (10.0.1.10)
COUNT   : 7 echecs / 45s   users: root, admin, backup
MITRE   : T1110.001 — Password Guessing
ACTION  : Ticket INC-2026-0412-001 cree — escalade N2
\end{verbatim}
\end{boxsiem}

% ==============================================================================
\section{Métriques de performance}
\label{sec:soc-kpi}
% ==============================================================================

\begin{boxdef}
\textbf{MTTD (\textit{Mean Time To Detect})~:} Durée moyenne entre la survenance
d'un incident et sa détection. Objectif niveau~4~: $<4$~heures incidents critiques.
\end{boxdef}

\begin{boxdef}
\textbf{MTTR (\textit{Mean Time To Respond})~:} Durée moyenne entre la détection
et les premières mesures de confinement. Objectif niveau~4~: $<4$~heures.
\end{boxdef}

\begin{boxdef}
\textbf{FPR (\textit{False Positive Rate})~:} Proportion des alertes ne
correspondant pas à un incident réel. Cible production~: $<30\,\%$.
\end{boxdef}

\begin{boxdef}
\textbf{Couverture ATT\&CK~:} Proportion des techniques MITRE ATT\&CK pour
lesquelles au moins une règle de détection est active dans le SIEM.
Exprimée par tactique (14~tactiques).
\end{boxdef}

\begin{table}[H]
\centering\small
\caption{Grille de maturité SOC (SOC-CMM adapté)}
\label{tab:soc-maturite}
\rowcolors{2}{TableauPair}{TableauImpair}
\begin{tabular}{p{1.2cm}p{2.5cm}p{3.8cm}p{3.5cm}p{2.2cm}}
\toprule
\tableauHeader{Niv.} & \tableauHeader{Désignation} &
\tableauHeader{Couverture} & \tableauHeader{Détection} &
\tableauHeader{MTTD cible} \\
\midrule
\textbf{1} & Initial   & Aucun SIEM, logs ad hoc       & Réactive (plainte user) & $>7$~j \\
\textbf{2} & Défini    & SIEM, règles par défaut        & Alertes non triées       & 2--7~j \\
\textbf{3} & Géré      & Règles personnalisées, triage  & FPR~$<50\,\%$            & $<24$~h \\
\textbf{4} & Optimisé  & Threat hunting, CTI intégrée  & FPR~$<30\,\%$            & $<4$~h \\
\textbf{5} & Adaptatif & SOAR, ML, automatisation       & Prédictive               & $<1$~h \\
\bottomrule
\end{tabular}
\end{table}

% ==============================================================================
\section{Threat Hunting et SOAR}
\label{sec:hunting-soar}
% ==============================================================================

\subsection{Threat Hunting}

\begin{boxdef}
\textbf{Threat Hunting~:} Processus proactif et itératif de recherche de signaux
de compromission latents, conduit par un analyste N3 à partir d'hypothèses
formulées via la CTI. Cherche ce que les règles SIEM ne voient pas encore.
\end{boxdef}

Cycle~: \textbf{Hypothèse} (CTI) $\to$ \textbf{Investigation} (requêtes SIEM/EDR)
$\to$ \textbf{Découverte} $\to$ \textbf{Amélioration} (nouvelle règle de détection).

\begin{lstlisting}[style=StyleTerminal,
  caption={Chasse aux PowerShell encodés --- requête KQL Kibana},
  label={lst:threat-hunt}]
# Hypothese : T1059.001 — PowerShell encode pour download de payload
# Source    : Sysmon Event ID 1 (Process Create)

event.code: "1"
  AND process.name: "powershell.exe"
  AND process.command_line: (*-EncodedCommand* OR *-enc* OR *-e* OR *-w hidden*)
  AND NOT process.parent.name: ("taskhostw.exe" OR "svchost.exe" OR "wmiprvse.exe")

# Connexions sortantes depuis PowerShell (Sysmon EID 3)
event.code: "3"
  AND process.name: "powershell.exe"
  AND NOT destination.ip: ("10.0.0.0/8" OR "192.168.0.0/16")
\end{lstlisting}

\subsection{SOAR}

\begin{boxdef}
\textbf{SOAR (\textit{Security Orchestration, Automation and Response})~:}
Plateforme automatisant les actions de réponse déclenchées par le SIEM~:
isolement réseau d'une machine, blocage IP au firewall, création de ticket,
notification. Exécute les playbooks sans intervention humaine pour les
actions de bas niveau, réduisant le MTTR.
\end{boxdef}

Distinction fondamentale~: le SIEM \textbf{détecte}, le SOAR \textbf{agit}.
Solutions open-source~: TheHive + Cortex, Shuffle. Commerciales~: Splunk SOAR,
Palo Alto XSOAR.

% ==============================================================================
\section*{Résumé}
% ==============================================================================

\begin{boxresume}
\begin{itemize}
  \item \textit{Assume breach}~: l'objectif est de détecter et contenir,
        pas seulement d'empêcher. Métrique clé~: le \textit{dwell time}.
  \item SOC~: capacité organisationnelle N1/N2/N3~; quatre modèles
        (interne, MSSP, hybride, vSOC)~; trois piliers (personnes,
        processus, technologie).
  \item SIEM~: pipeline en 6~étapes (collecte, parsing, normalisation,
        enrichissement, corrélation, alerte).
  \item Métriques~: MTTD, MTTR, FPR, couverture ATT\&CK.
  \item Threat hunting~: détection proactive par hypothèses CTI.
        SOAR~: automatisation de la réponse.
  \item LiZbiZ-SIS~: niveau~1 initial, objectif niveau~3.
\end{itemize}
\end{boxresume}

\begin{boxexo}
\textbf{Ex.~1.1 — Diagnostic de maturité.}
Appliquez la grille (tab.~\ref{tab:soc-maturite}) à LiZbiZ-SIS.
Justifiez chaque niveau et proposez deux actions correctives prioritaires.

\medskip\textbf{Ex.~1.2 — Sélection de sources.}
LiZbiZ-SIS dispose de~: 1~firewall pfSense, 2~serveurs Windows
(AD~+ ERP), 1~serveur Linux Apache, 1~MySQL, 30~postes Windows,
1~proxy Squid. Priorisez les 5~sources à intégrer en premier au SIEM
en justifiant par les menaces détectables.

\medskip\textbf{Ex.~1.3 — Règle de corrélation.}
Rédigez en syntaxe Wazuh ou Sigma une règle détectant un mouvement
latéral via SMB (T1021.002)~: connexion SMB réussie depuis un poste
utilisateur vers plus de 3~serveurs distincts en moins de 5~minutes.
Spécifiez niveau de gravité et mapping MITRE.
\end{boxexo}
